\section{배경}
\label{sec:bg}

\subsection{HAETAE 서명과 Fiat--Shamir with Aborts}
Shor 알고리즘~\cite{shor1994}에 의한 양자 위협으로 격자 기반 암호가 차세대 공개키 암호의 유력한
대안으로 부상하였으며~\cite{peikert2016}, 그 중 trapdoor 없는 Fiat--Shamir 격자
서명~\cite{lyubashevsky2012}이 표준화의 중심에 있다. Fiat--Shamir with
aborts~\cite{lyubashevsky2009}는 비밀키 $\svec$, 챌린지 $c$, 일회용 난수
$\yvec$로부터 응답 $\zvec=\yvec+c\svec$를 계산하고, $\zvec$가 비밀 정보를 노출할 위험이 있으면
해당 시도를 거부(abort)하고 새로운 난수로 서명을 재시도하는 거부 샘플링 구조를 사용한다.
\texttt{HAETAE}~\cite{haetae}는 여기에 bimodal 부호 $b\in\{0,1\}$로 $\pm c\svec$를 선택하고
초구 균등 분포에서 난수를 샘플링하여 서명 크기를 줄인다~\cite{bliss2013}. 다항식 연산은
$R_q=\mathbb{Z}_q[x]/(x^{256}+1)$, $q=64513$ 위에서 NTT로 수행되며, 본 논문은 보안 수준
120(MODE2: $L=4$, 비밀 다항식 $M=3$)을 기준으로 한다. 이때 거부 샘플링은 서명 경계 $B_1$(및
보조 bimodal 경계 $B_0$)을 사용하는 반면, 검증은 별도의 경계 $B_2$를 사용하며 \textbf{$B_0,B_1<B_2$}가
성립한다. 이 경계 차이는 \ref{sec:irv}장에서 단순 서명 후 검증의 한계를 설명하는 핵심 근거가 된다.

\subsection{오류주입공격 모델}
오류주입공격은 클럭/전압 글리치, 전자기파, 레이저 등을 통해 연산 중 디바이스의 상태를 일시적으로
교란하여, 출력된 오류 서명을 분석함으로써 비밀 정보를 복원하는 능동적 물리 공격이다. 오류 공격자 모델의 정형화는 이론과 실제를 잇는 관점에서 활발히 논의되어
왔다~\cite{toprakhisar2024,spruyt2021}. 본 논문은 \textbf{단일 오류} 모델을 가정하며, 그 효과를
(i)~명령어 스킵, (ii)~루프 조기 종료, (iii)~데이터 변조의 세 가지로 모델링한다. 이는 오류주입분석에서 널리 채택되는 표준 모델이다~\cite{espitau2016,bindel2016,ravi2022survey}. 공격자는 오류 상태에서 방출된 서명만
관찰할 수 있으며, 비밀키와 난수는 직접 관찰하지 못하고 챌린지 $c$ 등 공개 정보만을 이용한다. 일부
분석에서는 대응 기법의 검사 분기를 무력화하기 위한 \textbf{2차 오류}까지 고려한다.

\subsection{감염형 대응 기법의 종류}
오류 검출의 통상적 방어는 출력 직전에 결과를 검사하는 것이나, Yen과 Joye~\cite{yen2000}는 ``출력
전 검사''만으로는 오류 기반 분석을 충분히 막지 못함을 보였다. 감염형(infective) 대응 기법은 오류를
탐지한 뒤 연산을 중단하는 대신, 오류의 잔차를 출력에 확산시켜 결과 자체를 무효화한다. 즉 ``오류가
발생하면 중단한다''는 \emph{비교(분기) 연산}을 두지 않으므로, 그 분기를 건너뛰는 2차 오류에
본질적으로 강하다. Seo 등~\cite{seo2013}은
Fermat 소정리를 이용하여 비교 연산 없이 적은 추가 연산만으로 RSA 오류를 검출하는 기법을
제안하였으며, 본 논문은 이 설계 철학을 격자 기반 Fiat--Shamir 서명으로 확장한다.
